%!TEX root = ../main.tex
\section{Background and Setup}
\label{sec:bg}

\subsection{The URTIC corpus and the ComParE 2017 Cold sub-challenge}

The Upper Respiratory Tract Infection Corpus (URTIC) was distributed for the ComParE 2017 Cold sub-challenge~\cite{schuller2017interspeech}. It contains 28\,652 variable-length audio chunks (16\,kHz; median $5.2$\,s, mean $5.6$\,s, range $1.8$--$10.2$\,s, computed from the released waveform headers), drawn from $\sim$630 distinct speakers, with binary cold/non-cold (C/NC) labels at the chunk level. Class balance is heavily skewed towards non-cold ($\sim$9.5\,\% cold rate). The ``4students'' release we use exposes the train and devel subsplits but does not include speaker IDs; we cluster ECAPA-TDNN~\cite{desplanques2020ecapa} embeddings into $k{=}210$ pseudo-speakers (Section~\ref{ssec:pseudo}) as a proxy for the speaker probe protocol. The ComParE 2017 hidden test labels are not publicly available; we report all numbers on devel\_test (Section~\ref{ssec:splits}).

\subsection{Frozen WavLM-Large pooled-stats cache}

We use \texttt{microsoft/wavlm-large} (24 Transformer layers + 1 input layer, 1024-d hidden) without any fine-tuning. For each chunk, frame-level activations from layers $L \in \{0, 1, \dots, 24\}$ are mean+std+skew+kurt-pooled along the time axis after stripping the padding mask, producing a per-chunk fp16 tensor of shape $[25, 4096]$ ($25 \cdot 4 \cdot 1024 = 4 \cdot 25 \cdot 1024$ stats). The cache lives at \path|cache/microsoft_wavlm-large/pooled/{stem}.pt| and is reused by every downstream rung — A2, A2.5, A5b, A5.5, A6, A6b, A7, A7c, A7c-v2 all consume the same cache.

\subsection{Historical proxy-grouped sub-splits and overlap audit}
\label{ssec:splits}

The default URTIC distribution provides train/devel/test as three top-level partitions. Historically, we sub-split train into \texttt{train\_fit}/\texttt{train\_threshold} (90/10\,\%) and devel into \texttt{devel\_val}/\texttt{devel\_test} (50/50\,\%) with \texttt{StratifiedGroupKFold} keyed on Train-fitted $k{=}210$ pseudo-speaker assignments. This gives zero overlap under the same IDs used to construct the split. An independent side-local audit changes the interpretation: the Train split remains disjoint under recovered Train-local labels, but the canonical Development boundary is crossed by 201/210 Development-local groups, affecting 9,399/9,596 recordings; 27.46\,\% of recordings also have their nearest ECAPA neighbour across the boundary. These are proxy-connectivity diagnostics, not estimates or bounds of UAR optimism. A separate simplified-fusion comparison changed from $0.6317$ under the historical grouping to $0.6353$ under corrected pooled grouping, opposite to an optimism hypothesis and with high corrected-grouping variation. We therefore retract the speaker-disjoint wording but apply no numerical UAR correction. The historical Development results are pseudo-grouped partition diagnostics; load-bearing performance checks must use corrected side-local grouping or whole official sides.

\subsection{Pseudo-speaker clustering (k=210 ECAPA-TDNN)}
\label{ssec:pseudo}

\begin{figure*}[t]
\centering
\includegraphics[width=\textwidth]{fig1}
\caption{ECAPA-TDNN speaker-proxy diagnostics. \textbf{(A)} UMAP of the 19{,}101 train$+$devel ECAPA embeddings is descriptive only; clustering and validation use raw L2-normalised 192-d embeddings. \textbf{(B)} Tightness and separation of the historical Train-fitted $k{=}210$ map. \textbf{(C)} The pooled shipped-map comparison (all-recording ARI $0.309$) is a held-side fragmentation diagnostic, not validation of the shipped assignments. In side-local fits, all-recording KMeans--HDBSCAN ARI is $0.856/0.829/0.879$ on train/devel/test; excluding HDBSCAN noise gives $0.930/0.929/0.944$ at $97.28/96.69/97.59\,\%$ coverage. Thus ECAPA contains reproducible speaker-correlated structure, but its proxy IDs remain inferred and must be fitted within the side where they are used.}
\label{fig:ecapa}
\end{figure*}

\begin{table*}[t]
\centering\small
\caption{Reconciled raw-space ECAPA speaker-proxy diagnostics. KMeans stability
is the mean ARI over all 15 pairs from six seeds, not a single seed pair.
Non-noise ARI excludes HDBSCAN's unassigned points and is therefore reported
with coverage. The HDBSCAN cluster count is compared with the chosen $k=210$
prior, not a known true side-level speaker count. The final row deliberately
retains the disputed $0.309$ as a labeled held-side fragmentation diagnostic.}
\label{tab:ecapa_reconciliation}
\begin{tabular}{lrrrrrr}
\toprule
side & KMeans seed ARI & KM--HDB all & KM--HDB non-noise & coverage & HDB $k$ & shipped--local \\
\midrule
train & 0.889 & 0.856 & 0.930 & 97.28\% & 204 & 1.000 \\
devel & 0.897 & 0.829 & 0.929 & 96.69\% & 205 & 0.200 \\
test  & 0.888 & 0.879 & 0.944 & 97.59\% & 205 & 0.191 \\
\midrule
\multicolumn{7}{l}{Pooled Train+Development shipped-$k210$ versus pooled HDBSCAN, all recordings: ARI $=0.309$.} \\
\bottomrule
\end{tabular}
\end{table*}

Since URTIC's 4students release omits speaker IDs, the historical experiments cluster Train ECAPA-TDNN embeddings into $k{=}210$ pseudo-speakers and assign devel/test chunks by nearest Train centroid. Those assignments served as split keys, probe targets and masked-positive pools, but the side-local audit shows that the held-side assignments fragment Development and Test structure. As external calibration of the representation and clustering recipe, the same ECAPA--KMeans pipeline reaches ARI $0.963$ on 741 MLS-German recordings from 56 known speakers when fit at the known $k=56$; this does not supply URTIC ground truth. Corrected analyses fit proxy groups independently within each official side and prohibit cold labels from selecting the embedding, grouping method or $k$. These remain speaker-correlated proxies, not ground-truth identities.

\subsection{Speaker probe substrates}
\label{ssec:probes}

We use two probe substrates in matched form across every rung:

\textit{MLP probe} (\texttt{speakers.probe.train\_probe}): 2-layer MLP (256-d hidden, GELU, dropout 0.1) trained for 50 epochs at lr $10^{-3}$ on \texttt{(z, pseudo-speaker)} where $z$ is the 128-d projection output of the relevant rung. Training-time top-1 over $k{=}210$ classes; eval-time best top-1 reported.

\textit{LR probe} (\texttt{honesty.speaker\_probe}): scikit-learn multinomial logistic regression with C=1.0, max\_iter=2000. Operates on either the literal fusion-input vector (e.g., 3-d $[\ell_{\text{A2.5}}, z\ell_{g_4}, z\ell_{g_5}]$ for K=2 fusion) or the backbone-concat substrate (e.g., 4167-d \texttt{pooled\_4096} $\oplus$ \texttt{G4\_gi (7)} $\oplus$ \texttt{G5\_modulation (64)} for K=2 backbone-concat). The LR probe substrate is what we cite when speaker leakage is the load-bearing measurement axis.

\textit{Probe gate ceiling.} The acceptance threshold is $A2 \cdot \mathrm{LR\text{-}grouped} + 1\sigma = 0.0780$. Probes for any de-confounding rung must clear this ceiling on the matched substrate to count as ``not adding speaker leak.''

\textit{Substrate $\times$ split protocol matters} (M2 in our methodology contributions): the same A2 fused vector gives top-1 $= 0.0501$ under MLP-probe + grouped splits, $0.0760$ under LR-probe + grouped splits, $0.0501$ under MLP-probe + random splits, $0.0680$ under LR-probe + random splits. We anchor every gate on \texttt{honesty.speaker\_probe} substrate AND grouped splits — \emph{ceiling = $A2$ LR-grouped $+ 1\sigma = 0.0780$}.
